\documentclass{article}
\usepackage[dvips]{graphicx}
\usepackage{a4wide}
\usepackage{amsmath}
\usepackage{euscript}
\usepackage{amssymb}
\usepackage{amsthm}
\usepackage{amsopn}

\theoremstyle{definition}
\newtheorem*{definition}{Definition}
\newtheorem{theorem}{Theorem}

\renewcommand{\qed}{\blacksquare}

\newcommand{\vv}{\ensuremath{\vec{v}}}
\newcommand{\vu}{\ensuremath{\vec{u}}}
\newcommand{\vw}{\ensuremath{\vec{w}}}
\newcommand{\vx}{\ensuremath{\vec{x}}}
\newcommand{\vy}{\ensuremath{\vec{y}}}
\newcommand{\vb}{\ensuremath{\vec{b}}}
\newcommand{\vo}{\ensuremath{\vec{0}}}
\newcommand{\va}{\ensuremath{\vec{a}}}
\newcommand{\ve}{\ensuremath{\vec{e}}}

\newcommand{\R}{\mathbb{R}}
\newcommand{\Z}{\mathbb{Z}}
\newcommand{\C}{\mathbb{C}}
\newcommand{\N}{\mathbb{N}}
\newcommand{\Q}{\mathbb{Q}}

%%%%%%%%%%%%%%%%%%%%%%%%%%%%%%%%%%%%%%%%%%%%%%%%%%%%%%%%%%%

\begin{document}
	\begin{center}
		\Large{Math 251W: Foundations of Advanced Mathematics}

		\normalsize Portfolio problems from sections 3.2 $\&$ 3.3

		\hfill Name: Teddy Wachtler
		\hrule
		\vspace{0.4cm}
		\vspace{0.75cm}
	\end{center}

%%%%%%%%%%%%%%%%%%%%%%%%%%%%%%%%%%%%%%%%%%%%%%%%%%%%%%%%%%%

	\noindent Problem 3.3.17 [INDIVIDUAL LATEX]
	\vspace{0.3cm}

	Prove or give a counterexample

	\begin{enumerate}
		\item[1]
		Given sets $A$ and $B$, $\Pow (A\cup B)$ is necessarily equal to $\Pow (A)\cup \Pow(B)$.
		\vspace{.2cm}

		\fbox{Proof/Counterexample}

		\vspace{.15cm}

		We ask whether for all sets $A$ and $B$, $\Pow (A\cup B) = \Pow (A)\cup \Pow(B)$.
		We will show this to be false through a counterexample.
		Let $A$ and $B$ be particular sets so that $A = \{1\}$ and $B = \{2\}$.
		By the definition of set union, we see that $A \cup B = \{1, 2\}$.
		By the definition of the power set, we see that $\Pow (A\cup B) = \{\{ \emptyset\}, \{1\}, \{2\}\, \{1, 2\}\}$.
		By the definition of the power set, we see that $\Pow (A) = \{\{ \emptyset\}, \{1\}\}$ and $\Pow (B) = \{\{ \emptyset\}, \{2\}\}$.
		By the definition of set union, we see that $\Pow (A)\cup \Pow(B) = \{\{ \emptyset\}, \{1\}, \{2\}\}$.
		Since $\{1, 2\} \in \Pow (A\cup B)$ and $\{1, 2\} \not \in \Pow (A)\cup \Pow(B)$, we see $\Pow (A\cup B) \not = \Pow (A\cup B)$.
		Therefore, for all sets $A$ and $B$, $\Pow (A\cup B)$ is not necessarily equal to $\Pow (A)\cup \Pow(B)$.
		$\qed$.

		\vspace{.15cm}

		\item[2] \underline{proposition:}
		Given sets $A$ and $B$, $\Pow(A\cap B) = \Pow(A)\cap \Pow(B)$
		\vspace{.2cm}

		\fbox{Proof/Counterexample}

		\vspace{.15cm}

		We ask whether for all sets $A$ and $B$, $\Pow(A \cap B) = \Pow(A) \cap \Pow(B)$.
		Let $A$ and $B$ be arbitrary sets.
		Consider the following:

		Let $X$ be an arbitrary set in $\Pow(A \cap B)$.
		Let $x$ be an arbitrary element in $X$.
		By the definition of the power set, we see that $X \subseteq A \cap B$.
		By the definition of subsets, we see that $x \in A \cap B$.
		By the definition of set intersection, we see that $x \in A$ and $x \in B$.
		By the definition of subsets, since $x$ was an arbitrary element of $X$, we see that $X \subseteq A$ and $X \subseteq B$.
		By the definition of the power set, we see that $X \in \Pow(A)$ and $X \in \Pow(B)$.
		By the definition of set intersection, we see that $X \in \Pow(A) \cap \Pow(B)$.
		By the definition of subsets, since $X$ was an arbitrary element of $\Pow(A \cap B)$, we see that $\Pow(A \cap B) \subseteq \Pow(A)\cap \Pow(B)$,

		Let $Y$ be an arbitrary element in $\Pow(A) \cap \Pow(B)$.
		Let $y$ be an arbitrary element in $Y$.
		By the definition of set intersection, we see that $Y \in \Pow(A)$ and $Y \in \Pow(B)$.
		By the definition of the power set, we see that $Y \subseteq A$ and $Y \subseteq B$.
		By the definition of subsets, we see that $y \in A$ and $y \in B$.
		By the definition of set intersection, we see that $y \in A \cap B$.
		By the definition of subsets, since $y$ was an arbitrary element of $Y$, we see that $Y \subseteq A \cap B$.
		By the definition of the power set, we see that $Y \in \Pow(A \cap B)$.
		Since $Y$ was an arbitrary element of $\Pow(A) \cap \Pow(B)$, we see that $\Pow(A) \cap \Pow(B) \subseteq \Pow(A \cap B)$.

		Therefore, since $\Pow(A \cap B) \subseteq \Pow(A)\cap \Pow(B)$ and $\Pow(A) \cap \Pow(B) \subseteq \Pow(A \cap B)$, we see $\Pow(A \cap B) = \Pow(A) \cap \Pow(B)$.
		Therefore, for all sets $A$ and $B$, $\Pow(A \cap B) = \Pow(A) \cap \Pow(B)$.

		\vspace{.15cm}

	\end{enumerate}
	\vspace{1cm}

%%%%%%%%%%%%%%%%%%%%%%%%%%%%%%%%%%%%%%%%%%%%%%%%%%%%%%%%%%%

	\noindent Challenge Problems  \vspace{0.3cm}

	\begin{enumerate}
		\item[1]
		Given 3 arbitrary sets $A$, $B$, and $C,$ find an expression for $(A\cap B^C)\cup (B\cap C^C)$, using only union, intersection, and \emph{at most} one set difference.
		Prove the two expressions are equal.

		\item[2]
		Given a finite set $X$, with cardinality $|n|$ find $|\mathcal{P}(X\times X)\times\mathcal{P}(X\times X)|$ and $|\mathcal{P}(\mathcal{P}(X))|.$

	\end{enumerate}

%%%%%%%%%%%%%%%%%%%%%%%%%%%%%%%%%%%%%%%%%%%%%%%%%%%%%%%%%%%

\end{document}
